% Full-data rerun; its wide table is queued at the appendix entrance.
\section{Additional transfer and source-dependence analyses}
\label{app:fullscale}

\begin{table*}[t]
\centering
\caption{Source-cache repair at layer $16$ across four models and two benchmarks.
$n/G$ counts questions and skills retained using fresh baselines.
$S$, $C$, $SC$ and $SC+\mathrm{cut}$ show accuracy (\%).
The last three columns show cache benefit $\Delta_{\mathrm{cache}}$,
accuracy interaction $\Gamma$ and reading-cut loss $\Delta_{\mathrm{read}}$
in percentage points (Eq.~\eqref{eq:repair-effects}). Each row uses the
same questions for all conditions, with donor accuracy $100\%$ and
receiver accuracy $0\%$. These populations differ from Table~\ref{tab:paireddepth}.}
\label{tab:mechanism-contrasts}
\normalfont\rmfamily\fontsize{9}{11}\selectfont
\renewcommand{\arraystretch}{1.18}
\setlength{\tabcolsep}{2.5pt}
\begin{tabularx}{\textwidth}{@{}ll*{8}{>{\centering\arraybackslash}X}@{}}
\toprule
\rowcolor[HTML]{E4F4FA}
Model & Task & $n/G$ & $S$ & $C$ & $SC$ & $SC+\mathrm{cut}$ & $\Delta_{\mathrm{cache}}$ & $\Gamma$ & $\Delta_{\mathrm{read}}$ \\
\midrule
Qwen3-8B & MedCalc & $114/14$ & 3.5 & 7.0 & \textbf{42.1} & 7.0 & $+38.6$ & $+31.6$ & $+35.1$ \\
Qwen3-8B & TheoremQA & $53/43$ & 18.9 & 15.1 & \textbf{62.3} & 24.5 & $+43.4$ & $+28.3$ & $+37.7$ \\
\addlinespace[2pt]
Mistral-7B & MedCalc & $38/15$ & 7.9 & 18.4 & \textbf{18.4} & 5.3 & $+10.5$ & $-7.9$ & $+13.2$ \\
Mistral-7B & TheoremQA & $59/49$ & 3.4 & 5.1 & \textbf{10.2} & 6.8 & $+6.8$ & $+1.7$ & $+3.4$ \\
\addlinespace[2pt]
Llama-3.1-8B & MedCalc & $104/20$ & 3.8 & 5.8 & \textbf{18.3} & 9.6 & $+14.4$ & $+8.7$ & $+8.7$ \\
Llama-3.1-8B & TheoremQA & $72/56$ & 4.2 & 12.5 & \textbf{41.7} & 12.5 & $+37.5$ & $+25.0$ & $+29.2$ \\
\addlinespace[2pt]
DeepSeek-7B & MedCalc & $21/10$ & 0.0 & 19.0 & \textbf{28.6} & 19.0 & $+28.6$ & $+9.5$ & $+9.5$ \\
DeepSeek-7B & TheoremQA & $46/40$ & 2.2 & 6.5 & \textbf{10.9} & 8.7 & $+8.7$ & $+2.2$ & $+2.2$ \\
\bottomrule
\end{tabularx}
\end{table*}


The original Qwen3-8B/MedCalc controls sampled four rescued questions
from each of $40$ calculators ($160$ items). The group screen retained
$60$ from $15$ calculators. The rerun used all $467$ eligible questions
from $48$ calculators with the same interventions. Screening whole
groups now requires failure on up to $20$ questions instead of four,
retaining $135$ from $14$ calculators. Screening each item separately
retains $369$ from $45$ calculators (Table~\ref{tab:fullscale}).

For the dose and rank controls in Table~\ref{tab:fullscale}, the state
difference is $\delta^L=H_D^L[\mathcal I]-H_R^L[\mathcal I]$.
The rank implementation preserves the mean over source positions.
Let $\bar\delta^L$ repeat that mean at every row, and define
$\delta_k^L=\bar\delta^L+(\delta^L-\bar\delta^L)_k$, where
$(\cdot)_k$ retains the leading $k$ components of a singular value
decomposition, which ranks directions by how much matrix variation they capture.
The dose and rank controls share the notation
\begin{equation}
H_R^L[\mathcal I]\leftarrow H_R^L[\mathcal I]+\alpha\delta_k^L.
\label{eq:dose-rank}
\end{equation}
Full rank with $\alpha=1$ is full replacement, and $\alpha=0.5$ is half
the state difference. Rank controls use $\alpha=1$.
Increasing $k$ from $32$ to $64$ restores much of the recovery.
A bottom-$k$ control instead retains the least prominent directions of
the centered variation, with the same mean retained.
The second of four equal portions of the skill, containing its formula,
gives the largest recovery when copied alone in both samples.
The same-category mean donor changes from $0.05$ to $0.12$, on the $41$
full-rerun items where it is available. Here category refers to a calculator
group derived from skill-title keywords (for example, related fluid calculations). The intervention averages
donor-minus-receiver state differences from other calculators in that group.

\paragraph{The depth sweep at one-layer resolution.} The subset sweep measured
layers $0,4,8,12,14,15,16,20$ and placed the collapse somewhere after $12$.
At full scale we added $13$ and $17$ to $19$:
recovery is $0.98,0.98,0.89,0.74$ at layers $0,4,8,12$, then
$0.16,0.16,0.06,0.06,0.07,0.03,0.02,0.03$ at $13$ through $20$ under the group screen.
The item screen gives $1.00,0.98,0.94,0.79$ then
$0.32,0.34,0.25,0.25,0.23,0.19,0.14,0.17$.
Both screens show their largest local decline
between layers $12$ and $13$.

\paragraph{Broader screening and boundary sensitivity.}
The item-level screen retains the transfer decline. Paired knockout
comparisons use the common-item populations. Using transfer thresholds
$\tau=0.4$ or $0.6$ with dependence thresholds $1-\tau$ preserves the
boundary ordering in all eight main comparisons. Qwen3-8B/MedCalc dependence
first falls to at most half at layer $25$, rises above half again at $26$,
and remains below half from $27$ onward. The curves retain these recrossings.

\paragraph{Same-question joint failures.}
\label{app:joint-failure}
This check requires correct donors under both protocols, an incorrect
transfer receiver and an incorrect no-skill knockout baseline.
Joint failures at the transfer boundary are $100/124$ ($80.6\%$) for
Qwen3-8B at $T=13$ and $25/43$ ($58.1\%$) for Mistral-7B at $T=7$,
on $14$ and $15$ MedCalc skill groups. These sets differ from
Table~\ref{tab:paireddepth}. Four items in each pass the transfer
no-skill baseline despite failing it under knockout.

\paragraph{Repeated replacement as a consistency check.}
\label{app:repeated}
On the $135$ screened MedCalc items, replacing source states at layers
$8$ to $11$, $12$ to $15$, $14$ to $19$ and $16$ to $35$ recovers
$0.89,0.74,0.16,0.06$. Starting at the first layer recovers $0.98$.
Repeated and single writes produce identical tokens from each starting
layer, consistent with Appendix~\ref{app:cache-coverage}'s causal argument.

\paragraph{A continuous readout at the same boundary.}
We generate correct-skill reasoning once per item, then score that fixed
sequence under each intervention. This is teacher forcing: recorded
preceding tokens replace newly generated ones. Mean token log-probability
uses natural logarithms, measured in nats. Higher values indicate stronger
predictions of the recorded sequence. We normalize it as
$z=(\ell_L-\ell_{\mathrm{recv}})/(\ell_{\mathrm{gold}}-\ell_{\mathrm{recv}})$,
where $\ell$ is the mean log-probability for the patched, receiver or
donor condition. This $z$ differs from knockout retention $q$.

On $442$ donor-correct items from $48$ calculators, $z$ is
$0.999,0.989,0.960,0.831,0.598,0.595,0.340,0.201,0.129$ at layers
$0,4,8,12,13,14,16,20,24$. The $0.233$ drop from $12$ to $13$ is near
the accuracy boundary in the separate $134$-item depth set.
The decline spans all four quarters of the reasoning sequence
($0.20,0.26,0.31,0.21$). After supplying the reasoning, all conditions
predict the final number closely: median log-probabilities lie between
$-0.001$ and $-0.005$. Their donor-minus-receiver denominator has median
$0.0014$ and is below $0.01$ on $59\%$ of items, limiting this readout.

\paragraph{Distributed reading and receiver state.}
The attention output contributed by skill positions,
normalized by the length of the hidden-state vector, concentrates in
layers $12$ to $21$. Blocking the $32$ highest-ranked
heads costs $0.39$ of the skill's effect, versus $0.23$ for random heads at the
same layers (calculators not used to rank the heads, paired difference $0.154$). Reading uses many heads rather than a small subset. After late span replacement, the attention output from skill positions remains
$0.86$ to $0.99$ of the donor's despite weak behavioural recovery. Earlier receiver question-state updates remain unchanged, motivating
the interaction tests in \S\ref{sec:mechanism}.


\subsection{Source-cache repair and attention-path experiments}
\label{app:mechanism}

The focused Qwen3-8B/MedCalc experiments fix $\mathcal Q_0$, comprising
$124$ historical IDs from $14$ calculators, before inference. Four NVIDIA
RTX 4090 graphics processors each receive every fourth ID, using the
$1{,}536$-token protocol in Appendix~\ref{app:detail}.
Let $e_i^X$ be binary correctness for item $i$ in condition $X$, with
$0$ denoting no skill. Selection uses only fresh baselines:
\begin{equation}
\mathcal Q_{\mathrm{rep}}
=\{i\in\mathcal Q_0:e_i^D=1,\ e_i^R=e_i^0=0\}.
\label{eq:repair-set}
\end{equation}
Of $116$ donor-correct, receiver-incorrect items, $114$ also fail without
a skill. On these same $114$ items, every condition uses
$a_X(L)=|\mathcal Q_{\mathrm{rep}}|^{-1}
\sum_{i\in\mathcal Q_{\mathrm{rep}}}e_i^X(L)$.
All $372$ baseline generations agree
token-for-token between the two suites. The complete repair grid contains
$2{,}356$ generations, the path grid $1{,}488$, and worker identity checks
$100$. All $48$ own-state writebacks are token-exact.

Figure~\ref{fig:mechanism} and Table~\ref{tab:mechanism-contrasts} use this
rule for eight comparisons with separate candidates. Repair and path tests
share item IDs and baseline tokens within each row. Repair is measured
at $L8$, $L12$, $L16$ and $L24$, and path cuts at $L8$, $L16$ and $L24$.
Writes, masks and contrasts follow
Eqs.~\eqref{eq:cache}, \eqref{eq:reading-cut} and \eqref{eq:repair-effects}.

\paragraph{Qwen3-8B/MedCalc truncation sensitivity.}
We retain all intervention outcomes, including outputs stopped at the
generation limit, which we call truncated outputs. For the
prespecified difference between $SC$ and $S$ at layer $16$, the gain is $38.6$ percentage
points.
The corresponding accuracy interaction $\Gamma$ is $31.6$ points. $SC-(SC+\mathrm{cut})$ is $35.1$ points.
On the $96$ items with all three baselines untruncated, the $SC$ gain over $S$ is
$38.5$ points.
At layer $16$, $S$, $SC$ and $SC+\mathrm{cut}$ truncate $4$, $4$ and $3$ of $114$ outputs.
Scoring truncated $S$ answers as correct and truncated $SC$ answers as
incorrect still leaves a
$35.1$-point repair advantage. At layer $24$, $SC$ and $SC+\mathrm{cut}$ truncate $3$ and $28$
outputs and show no additional loss in measured accuracy.

At layer $24$, reader and source-self cuts retain $29/114$ and $110/114$
correct answers. Their loss difference is $71.1$ points.
Reader cuts truncate $21$ outputs. Counting all as correct still leaves
retention at $50/114$, below source-self retention $110/114$ (no truncation).
Both contrasts survive this worst-case reassignment. Experiment records retain counts for each condition, file checksums
and these sensitivity calculations.

\paragraph{Redundant combined cut.}
With causal attention and otherwise positionwise operations, earlier
tokens cannot read the skill. Persistent reading cuts also prevent its
changed states from reaching later tokens. Adding a within-skill cut
therefore leaves outputs unchanged by induction over layers and generated
tokens. We omit this redundant condition. The argument does not apply
when later reading is allowed.

\subsection{Cache coverage under causal prefill}
\label{app:cache-coverage}

This deduction assumes a deterministic causal Transformer with masked
attention as its only cross-position operation
\citep{vaswani2017attention}, identical parameters and position IDs,
and the same exact-arithmetic computation without extra attention cuts.
Donor and receiver tokens differ only within the complete, contiguous
source interval $\mathcal I=[\mathrm{lo},\mathrm{hi})$ in the prefix.

\paragraph{Proposition (source-cache coverage).}
Suppose $S$ replaces the entire source output of block $L$ before the
next block executes, and the cache write occurs only after prefix prefill.
At the observation point in Eq.~\eqref{eq:continuation},
\begin{equation}
\mathcal C_S^\ell[\mathcal I]=
\begin{cases}
\mathcal C_R^\ell[\mathcal I],&\ell\le L,\\
\mathcal C_D^\ell[\mathcal I],&\ell>L.
\end{cases}
\label{eq:source-coverage}
\end{equation}
Together with the explicit write at layers $0$ through $L$, this gives
$\mathcal C_{SC}[\mathcal I]=\mathcal C_D[\mathcal I]$.

\emph{Proof.}
Block-$L$ KV is computed before the output patch, leaving source caches
at $\ell\le L$ unchanged. Before $\mathrm{lo}$, identical tokens and
causality give identical states at every layer. After $S$, all block-$L$
outputs at $i<\mathrm{hi}$ therefore match the donor.
At these positions, the next block reads only matching earlier or
current states, producing matching source KV and outputs.
Induction proves the upper-layer case.
Equation~\eqref{eq:cache} supplies the lower layers for $SC$.
\hfill$\square$

\begin{table}[t]
\centering
\caption{Prefix-cache coverage before the final prompt token, assuming
deterministic causal computation in exact arithmetic and a complete,
contiguous source span with aligned positions and identical preceding
tokens. $D$ and $R$
denote donor and receiver source KV. $\mathcal B_R$ and $\mathcal B_S$
are all other prefix caches from the unpatched and state-patched runs.
The entries are identities derived under these assumptions.}
\label{tab:cache-coverage}
\small
\renewcommand{\arraystretch}{1.18}
\setlength{\tabcolsep}{5pt}
\begin{tabular}{@{}lccc@{}}
\toprule
Condition & \multicolumn{2}{c}{Source KV} & Other KV \\
\cmidrule(lr){2-3}
 & $\ell\le L$ & $\ell>L$ & All layers \\
\midrule
Receiver & $R$ & $R$ & $\mathcal B_R$ \\
$S$ & $R$ & $D$ & $\mathcal B_S$ \\
$C$ & $D$ & $R$ & $\mathcal B_R$ \\
$SC$ & $D$ & $D$ & $\mathcal B_S$ \\
\bottomrule
\end{tabular}
\end{table}


\paragraph{Remaining execution state.}
The post-prefill write preserves the other caches
(Table~\ref{tab:cache-coverage}). Under these assumptions,
\begin{equation}
\begin{aligned}
\hat y_{SC}&=\mathcal G_{\theta,M_0}
 (u_*,\mathcal C_D[\mathcal I],\mathcal B_S),\\
\hat y_D&=\mathcal G_{\theta,M_0}
 (u_*,\mathcal C_D[\mathcal I],\mathcal B_D),
\end{aligned}
\label{eq:source-not-full}
\end{equation}
where $M_0$ is the original mask. Since $\mathcal B_S$ need not equal
$\mathcal B_D$, answers may differ. All-layer source-KV-only replacement
instead preserves $\mathcal B_R$ and is generally not equivalent to $SC$.

Under $SC+\mathrm{cut}$, source queries remain unmasked, so the same
induction applies. However, $\mathcal B_{SC+\mathrm{cut}}$ may differ
from $\mathcal B_S$, and continuation uses a different mask. The contrast
includes both prefill and decoding effects.

Code inspection confirms the stated write ordering. Exact source-cache
identities rely on the arithmetic assumption above. Discrepancies against
the donor in BF16 runs remain unmeasured.

\subsection{A sufficient condition for behavioral recovery}
\label{app:behavior-recovery}

We apply margin-based stability reasoning \citep{tsuzuku2018lipschitz}
to cached decoding, adding a sensitivity assumption to Eq.~\eqref{eq:cache-distance}.

\paragraph{One-step guarantee.}
Fix the model, last prompt token $u_*$, positions and original mask $M_0$.
For vocabulary $\mathcal V$, let $F(\mathbf z)\in\mathbb R^{|\mathcal V|}$
map initial prefix caches to next-token logits. Assume finite sensitivity
$\kappa$ over the compared states in $\Omega$. The infinity norm is the
largest absolute change in any token score:
\begin{equation}
\|F(\mathbf z)-F(\mathbf z')\|_\infty
\le\kappa\|\mathbf z-\mathbf z'\|_2,\qquad
\mathbf z,\mathbf z'\in\Omega.
\label{eq:cache-lipschitz}
\end{equation}
For the donor's unique greedy choice $y_D$, define its positive margin
\begin{equation}
m_D=F_{y_D}(\mathbf z_D)-\max_{v\ne y_D}F_v(\mathbf z_D)>0.
\label{eq:donor-margin}
\end{equation}
\emph{Proposition.} If $2\kappa\varepsilon_{SC}<m_D$, then greedy
decoding from $\mathbf z_{SC}$ selects $y_D$.
\emph{Proof.} Each logit changes by at most
$\kappa\varepsilon_{SC}$. Thus, for every $v\ne y_D$,
\begin{equation}
F_{y_D}(\mathbf z_{SC})-F_v(\mathbf z_{SC})
\ge m_D-2\kappa\varepsilon_{SC}>0.
\label{eq:margin-preservation}
\end{equation}
The factor $2$ allows the winning logit to decrease and a competitor
to increase. Strict inequality excludes ties. \hfill$\square$

\paragraph{Complete generation.}
Let the donor produce $N$ tokens including its stopping token.
$F_t^D(\mathbf z)$ processes $u_*$ and fixed donor history $y_{D,<t}$
from initial cache $\mathbf z$, recomputing all new caches.
Assume Eq.~\eqref{eq:cache-lipschitz} for $F_t^D$ with bound $\kappa_t$
and donor margin $m_{D,t}>0$. The bound includes all intervening cache updates:
\begin{equation}
\bigl[\ \forall t\le N:\ 2\kappa_t\varepsilon_{SC}<m_{D,t}\ \bigr]
\quad\Longrightarrow\quad \hat y_{SC}=\hat y_D.
\label{eq:sequence-recovery}
\end{equation}
The first token agrees by the one-step result. If earlier tokens agree,
decoding follows the history defining $F_t^D$, so token $t$ agrees.
Induction under the common stopping rule gives an identical sequence,
which is correct whenever the donor's answer is correct.

\paragraph{Scope.}
Common valid $\kappa_t$ over $D$, $S$ and $SC$ give $SC$ no larger error
bounds (Eq.~\eqref{eq:cache-distance}). Losing cancellation can increase
actual errors, so accuracy gains are not guaranteed. Different sequences
may answer correctly, making this condition only sufficient.
$SC+\mathrm{cut}$ changes the mask and prefix and is excluded.
Sensitivity bounds and itemwise margins remain unverified.


